\chapter{Orthopedics}

\medterm{Achilles Tendonitis}-- Inflammation and degeneration of the Achilles tendon, most common in runners and middle-aged athletes. Insertional: at the calcaneal insertion, associated with bone spurs and Haglund deformity. Non-insertional (mid-portion): 2-6 cm above insertion, due to repetitive microtrauma and tendinosis. Presents with posterior heel pain and morning stiffness, improving with activity and worsening with activity (later stages). Diagnosis: clinical; ultrasound/MRI for confirmation. Treatment: eccentric calf stretching (most effective), activity modification, heel lifts, NSAIDs, and shockwave therapy. Refractory cases: PRP, tenotomy, or surgical debridement.

\medterm{Achilles Tendon Rupture}-- Complete tear of the Achilles tendon, most common in middle-aged men during recreational sports (weekend warrior). Mechanism: sudden push-off with knee extended. Presents with hearing/feeling a pop, acute posterior ankle pain, inability to stand on tiptoes, and palpable gap in the tendon. Positive Thompson test (no plantarflexion with calf squeeze). Diagnosis: clinical; MRI/ultrasound for confirmation. Treatment: non-operative (functional bracing in equinus with early weight-bearing, complication: rerupture 5-10\%) vs. operative repair (open vs. percutaneous, complication: wound healing, infection 2-5\%). Both have similar long-term outcomes.

\medterm{ACL Injuries}-- Anterior cruciate ligament (ACL) tears: non-contact (pivot, deceleration) > contact. Clinical: Lachman (most sensitive), pivot shift (most specific), anterior drawer. MRI confirms (bone bruise lateral femoral condyle + lateral tibial plateau, ACL discontinuity). Treatment: nonoperative (rehab, bracing) for sedentary, low-demand; ACL reconstruction (autograft: BTB - gold standard, hamstring, quadriceps tendon; allograft: higher failure in young active) for active, instability. Timing: acute (<3 wacute (<3 weeks): nonoperative vs operative (controversial, no difference in outcomes); subacute to chronic: reconstruction. Post-op: brace locked 0-90, NWB 6 weeks, WBAT 6-12 weeks, return to sports 6-9 months.

\medterm{ACL Reconstruction}-- Surgical reconstruction of the anterior cruciate ligament (ACL) following rupture, one of the most common knee ligament injuries. The ACL is a primary restraint to anterior tibial translation and rotatory instability, and its rupture frequently occurs during non-contact pivoting sports activities ( skiing, soccer, basketball). Risk factors include female sex (2-8 times higher incidence than males, attributed to hormonal, anatomic, and neuromuscular factors), family history, participation in highhigh-risk sports with cutting, pivoting, and jumping. The surgical technique involves harvesting an autograft (bone-patellar tendon-bone or hamstring tendon graft are the most common) or using an allograft for multi-ligament knee injuries.

\medterm{Acromioclavicular Joint Injuries}-- AC joint injuries: Rockwood classification: I (sprain AC ligaments), II (AC ligaments torn, CC intact), III (AC+CC torn, 25-100\% displacement), IV (posterior displacement into trapezius), V (superior displacement >100\%), VI (inferior displacement subcoracoid/subacromial). Treatment: I/II: nonoperative (sling, PT); III: controversial (nonoperative for most, surgical for high-demand overhead athletes, laborers); IV-VI: surgical (anatomic CC reconstruction: allograft/autograft, synthetic tape, sutusutures, cerclage wires). Complications: AC joint arthritis, distal clavicle osteolysis, coracoclavicular ossification, fixation failure, and infection.

\medterm{Adhesive Capsulitis}-- A condition of progressive shoulder stiffness and pain caused by inflammation and fibrosis of the glenohumeral joint capsule. Common name: frozen shoulder. Stages: freezing (painful, 2-9 months), frozen (stiff, 4-12 months), thawing (gradual improvement, 5-24 months). Most common ages 40-60. Risk factors: diabetes (strongest), thyroid disease, female sex, and prior shoulder trauma/surgery. Bilateral in up to 40\% of diabetics. Diagnosis: clinical (restricted active and passive ROM in all planes, especially external rotation). Treatment: physical therapy, NSAIDs, intra-articular corticosteroid injection, hydrodilatation (capsular distension), and arthroscopic capsular release for refractory cases.

\medterm{Adult Spinal Deformity}-- Adult spinal deformity (ASD): scoliosis (coronal Cobb >10 deg), kyphosis (sagittal), sagittal imbalance. SRS-Schwab classification: PI-LL mismatch, SVA (sagittal vertical axis), PT (pelvic tilt). Clinical: back pain, radiculopathy, neurogenic claudication, disability (ODI, SRS-22). Imaging: full-length standing AP/lateral, bending, CT for bone quality. Treatment: nonoperative (PT, core, epidural injections); operative: decompression + fusion (PSF, ALIF, TLIF, XLIF, OLIF, anterior column realignmanterior column realignment, three-column osteotomy, PSO). Complications: pseudarthrosis, implant failure, proximal junctional kyphosis (PJK), infection, neurologic deficit. Outcomes: reliable pain relief, functional improvement, but high complication rates (30-40 percent).

\medterm{Amelia}-- A rare congenital anomaly characterized by complete absence of one or more limbs. Unilateral upper limb amelia is the most common form. Causes include vascular disruption, amnion rupture sequence, and teratogens. May be isolated or syndromic (Roberts syndrome, ADAM sequence, amniotic band syndrome). Diagnosis by prenatal ultrasound. Management: prosthetic fitting as early as 6 months of age to promote functional adaptation and normal development. Multidisciplinary care includes physical and occupational therapy. Children with amelia can achieve high levels of independence with adaptive technology.

\medterm{Ankle Arthroplasty}-- Total ankle arthroplasty (TAA): indications: end-stage ankle OA (post-traumatic, primary, inflammatory), preserved hindfoot alignment (<10-15 deg varus/valgus), no severe osteoporosis, no active infection. Contraindications: severe deformity, Charcot, prior infection, young high-demand. Implants: 3-component (mobile bearing: STAR, Salto, Hintegra) vs 2-component (fixed bearing: Infinity, Vantage). Approaches: anterior (standard), anterolateral. Alignment: neutral mechanical axis. Complications: Complications: infection, wound healing problems, deep vein thrombosis, aseptic loosening, impingement, gutter pain, subsidence, and polyethylene wear.

\medterm{Ankle Fractures}-- Ankle fractures: Weber classification (fibular fracture relative to syndesmosis): A (below syndesmosis, stable), B (at syndesmosis, variable stability), C (above syndesmosis, unstable). Lauge-Hansen: supination-external rotation (SER, most common), supination-adduction, pronation-external rotation, pronation-abduction. Treatment: stable (Weber A, SER stage I/II): closed reduction, short leg cast/boot, NWB 6 weeks; unstable (Weber B/C, SER III/IV): ORIF (lateral plate for fibula, medial screws/temedial screws/tension band for medial malleolus), syndesmotic screw fixation (if unstable, 3.5-4.5 mm quadricortical screw 2-3 cm above tibiotalar joint, weightbearing at 8 weeks with screw removal at 3-4 months; alternatively, suture button fixation). Post-op: splint, NWB 6-8 weeks, then progressive weightbearing. Complications: wound breakdown, infection, malunion (most common), nonunion, post-traumatic osteoarthritis (most common late complication).

\medterm{Ankle Sprain}-- An injury to the ligaments surrounding the ankle joint, most commonly involving the lateral ligament complex (anterior talofibular ligament [ATFL, most commonly injured], calcaneofibular ligament [CFL], and posterior talofibular ligament [PTFL]) from an inversion mechanism. Medial (deltoid) ligament sprains occur less commonly from eversion injuries and syndesmotic (high ankle) sprains from external rotation. Classification: Grade I (mild, microscopic ligament fibers, minimal sweswelling and ecchymosis with a stable ankle joint), Grade II (moderate, partial tearing of ligament fibers, with moderate swelling, ecchymosis, and joint laxity with a firm end point), and Grade III (severe, complete ligament rupture, with profound swelling, ecchymosis, and joint instability with no firm end point).

\medterm{Ankylosing Spondylitis}-- A chronic, progressive inflammatory disease primarily affecting the axial skeleton (sacroiliac joints and spine), characterized by inflammation at entheses (where tendons, ligaments, and joint capsules insert into bone), leading to new bone formation, syndesmophyte development, and eventual vertebral fusion (bamboo spine). It is the prototype of the spondyloarthritis family of diseases and is strongly associated with HLA-B27 (present in 80-90 percent of patients, though HLA-B27 is present in 6-86-8 percent of the general population, with the risk of developing AS in HLA-B27-positive individuals being 1 to 2 percent per year. The disease typically begins in early adulthood with insidious onset of low back pain and stiffness that is worse in the morning and with rest, relieved by activity.

\medterm{Anterior Cruciate Ligament Injury}-- Tear of the ACL, a common sports injury (non-contact deceleration, pivoting, or hyperextension). Presents with a pop, immediate swelling (hemarthrosis), and giving way of the knee. Positive Lachman test (most sensitive), anterior drawer test, and pivot shift test. Associated injuries: MCL (O'Donoghue triad: ACL, MCL, medial meniscus), meniscal tears. Diagnosis: MRI. Treatment: functional rehab for low-demand patients; ACL reconstruction (hamstring autograft, patellar tendon autograft, quadriceps tendon, or allograft) for high-demand athletes or those with persistent instability.

\medterm{Arthrogryposis}-- A group of congenital disorders characterized by multiple joint contractures present at birth, caused by decreased fetal movement (fetal akinesia) during development. Etiology includes neuropathic (most common, spinal anterior horn cell degeneration), myopathic (congenital muscular dystrophies), and mechanical (oligohydramnios, uterine abnormalities). Amyoplasia is the most common type (1 in 3,000 births): symmetric limb involvement, no internal organ involvement, and normal cognition. Physical exam shows fixed joint contractures (clubfeet, flexed wrists, extended knees), smooth, shiny skin over affected joints, and decreased muscle bulk. Joints are typically rigid with no skin creases. Treatment: aggressive physical therapy, serial casting, and orthopedic surgery (tendon releases, osteotomies) to improve function. Most children achieve ambulation with assistive devices. Intelligence is typically normal.

\medterm{Articular Cartilage Injuries}-- Chondral/osteochondral lesions: Outerbridge grading (I: softening, II: fibrillation <50\%, III: fibrillation >50\%, IV: subchondral bone exposed). OCD (osteochondritis dissecans): subchondral bone necrosis, fragment separation (knee: lateral femoral condyle, elbow: capitellum, ankle: talar dome). Treatment: microfracture (marrow stimulation, fibrocartilage, <2-3cm2), OATS (osteochondral autograft transfer, mosaicplasty), allograft transplantation (large defects), MACI (matrix-induced autologous chchondrocyte implantation). Treatment: repair if possible (young, active, large fragment) vs excision (small, loose).

\medterm{Backache}-- Pain or discomfort located anywhere in the posterior trunk from neck to buttocks, a symptom rather than a disease. Causes include musculoskeletal strain (most common), disc herniation, spinal stenosis, facet arthropathy, sacroiliac joint dysfunction, and referred pain from renal, pancreatic, or vascular sources. Acute backache often resolves with conservative care; chronic backache benefits from exercise, physical therapy, and multidisciplinary pain management.

\medterm{Back Pain}-- One of the most common medical complaints, affecting up to 80 percent of the population at some point in life, and the leading cause of disability worldwide. Back pain is classified by duration: acute (less than 4 weeks), subacute (4-12 weeks), and chronic (greater than 12 weeks). It is further divided into mechanical (non-specific) back pain (90 percent of cases, caused by muscle strain, ligament sprain, degenerative disc disease, facet joint arthropathy), radicular pain (from nerve root comprecompression (from herniated disc, spinal stenosis, or foraminal stenosis), and referred pain (from conditions such as pancreatitis, renal colic, aortic aneurysm). The diagnosis is clinical, based on history and physical examination. Imaging is indicated when red flags are present.

\medterm{Baker Cyst}-- A fluid-filled cyst in the popliteal fossa, resulting from herniation of the synovial membrane through the posterior knee joint capsule or from a communication between the knee joint and the gastrocnemius-semimembranosus bursa. A Baker cyst is not a true cyst but rather an effusion of the gastrocnemius-semimembranosus bursa that communicates with the knee joint, containing synovial fluid. It occurs in association with knee joint effusion from any underlying cause, most commonly osteoarthritis, rrheumatoid arthritis, or trauma, and it presents as a palpable, firm, mildly tender mass in the popliteal fossa that is more prominent with knee extension (Foucher sign). Diagnosis is confirmed by ultrasound or MRI.

\medterm{Bone Metastases}-- Bone metastases: most common from breast, prostate, lung, kidney, thyroid. Solitary vs multiple. Clinical: pain (worse at night), pathologic fracture, hypercalcemia, spinal cord compression. Mirels score for impending pathologic fracture: site (upper limb 1, lower limb 2, peritrochanteric 3), pain (mild 1, moderate 2, functional 3), lesion (blastic 1, mixed 2, lytic 3), size (<1/3 1, 1/3-2/3 2, >2/3 3). Score >=9: prophylactic fixation. Spinal metastases: Bilsky epidural spinal cord compression compression (Bilsky grade 0-3). Grade 1: external beam radiotherapy + bisphosphonates/denosumab; grade 2/3: surgical decompression plus stabilisation followed by radiotherapy. Pathologic fracture: surgical fixation or arthroplasty.

\medterm{Bones and Joints}-- Bones and joints form the skeletal system, providing structural support, protection of organs, and enabling movement through articulations between bones. Joints are classified as fibrous, cartilaginous, or synovial, with synovial joints offering the greatest range of motion. Common conditions include osteoarthritis, rheumatoid arthritis, fractures, and osteoporosis.

\medterm{Brachial Plexus Injuries}-- Brachial plexus injuries: traction (motorcycle accidents, sports), compression (crutches, thoracic outlet). Classification: supraclavicular (roots/trunks), infraclavicular (divisions/cords). Narakas classification: I (proximal C5-C6, Erb's palsy: waiter's tip), II (C5-C7), III (C5-T1), IV (complete C5-T1), V (complete + root avulsions). Root avulsion: pseudomeningocele on MRI myelography, Horner syndrome (T1 avulsion), SER test (sensory evoked potentials). Treatment: early (3-6 months): neurolysneurolysis vs reconstruction (sural nerve graft, nerve transfer). Late (>6 months): tendon transfer, free muscle transfer, arthrodesis. Prognosis: better for lower trunk/infraclavicular, worse for upper trunk/supraclavicular.

\medterm{Bunion}-- A bony prominence (exostosis) at the medial aspect of the first metatarsal head, causing lateral deviation of the great toe (hallux valgus). The deformity is characterized by a structural abnormality at the first metatarsophalangeal joint, with the first metatarsal deviating medially and the proximal phalanx deviating laterally. The bunion itself is the medial eminence of the first metatarsal head, which becomes prominent and painful due to chronic pressure from footwear. Risk factors include nanarrow-toed shoes, high heels, female sex (10:1), family history, and inflammatory arthritis. Diagnosis is made on physical examination with weight-bearing radiographs confirming the hallux valgus angle and intermetatarsal angle.

\medterm{Bursitis}-- Inflammation of a bursa, a fluid-filled sac that reduces friction between tendons, muscles, bones, and skin. Bursae are lined by synovial cells and normally contain a thin layer of lubricating synovial fluid. Bursitis most commonly affects the subacromial (shoulder), olecranon (elbow), trochanteric (hip), prepatellar (knee), infrapatellar, and retrocalcaneal (ankle) bursae. Causes include repetitive microtrauma or overuse (the most common mechanism), acute trauma, infection (septic bursitis, moscommonly Staphylococcus aureus (80 to 90 percent of cases), which may complicate superficial skin infections, rheumatoid arthritis, and sporting injuries, with the olecranon and prepatellar bursae being the most common sites.

\medterm{Carpal Tunnel Syndrome}-- See neurology.

\medterm{Cervical Spine Trauma}-- Cervical spine trauma: Subaxial (C3-C7) injury classification (SLIC - Subaxial Cervical Spine Injury Classification): morphology (1-4), disco-ligamentous complex (0-2), neurologic (0-3). Total <4 nonop, >4 op, =4 surgeon discretion. Upper cervical: occiput-C1 (occipital condyle fracture, occipitocervical dissociation), C1 (Jefferson burst fracture: lateral mass displacement >6.9mm on CT), C2 (odontoid fracture: Type I tip, Type II base - most common, nonunion risk >5mm displacement/age >70/smokismoking/diabetes); Type III (odontoid body/synchrondrosis, less than 3 percent, surgical). Treatment: nonoperative (collar, Halo vest) for stable; C1 lateral mass screws (Harms), C1-2 transarticular screws (Magerl), occipitocervical fusion for unstable.

\medterm{Chondrosarcoma}-- The second most common primary malignant bone tumor (after osteosarcoma) and the most common primary malignant bone tumor in adults (peak age 40-70 years). It arises from cartilage-forming cells and is classified by grade (Grade 1: low-grade/conventional; Grade 2: intermediate; Grade 3: high-grade) and by location (central/intramedullary, surface/peripheral, or secondary). Conventional chondrosarcoma typically arises in the proximal femur, pelvis, and proximal humerus. Secondarary chondrosarcoma may arise from a pre-existing osteochondroma (in multiple hereditary exostosis when the cartilage cap thickness exceeds 1.5 to 2 cm, the risk of malignant transformation is approximately 2 to 5 percent) or from Paget disease of bone. The clinical presentation is with localised pain, a palpable mass, and pathological fracture in approximately 10 to 15 percent of cases.

\medterm{Clavicle Fractures}-- Clavicle fractures: 80\% midshaft (Group I), 15\% lateral (Group II: Neer I nondisplaced, II displaced with CC ligament disruption, III intra-articular), 5\% medial (Group III). Midshaft: nondisplaced/shortening <2cm: sling, figure-8; displaced/comminuted/shortening >2cm/skin tenting: ORIF (superior plate, anteroinferior plate). Lateral: Neer I: nonoperative; Neer II: ORIF (hook plate, anatomic plate, CC screw + plate); Neur III: ORIF. Medial: rare, check for mediastinal injury. Complications: nonunonunion (5-15 percent), malunion, hardware irritation, shoulder stiffness.

\medterm{Clubfoot (Talipes Equinovarus)}-- A congenital foot deformity with an incidence of approximately 1 in 1,000 live births, characterized by four components: equinus (ankle plantarflexion), varus (hindfoot inversion), adductus (forefoot adduction), and cavus (high arch). Clubfoot is classified as idiopathic (most common, no identifiable cause), postural (positional, passively correctable), or syndromic (associated with neuromuscular conditions like myelomeningocele, arthrogryposis, or chromosomal abnormalities). The male-to-female male-to-female ratio is 2:1. The Ponseti method is the gold standard for treatment, involving serial casting, Achilles tenotomy, and bracing.

\medterm{Codman Triangle}-- A characteristic periosteal reaction seen on plain radiographs, formed when a aggressive bone lesion elevates the periosteum away from the bone cortex, creating a triangular area of new reactive bone at the angle between the raised periosteum and underlying cortex. Highly suggestive of malignant bone tumors (most commonly Osteosarcoma and Ewing Sarcoma) or severe osteomyelitis.

\medterm{Compartment Syndrome}-- A surgical emergency caused by increased pressure within a closed osteofascial compartment, leading to tissue ischemia, necrosis, and permanent functional loss. The lower leg (anterior, lateral, deep posterior, and superficial posterior compartments) and forearm (flexor and extensor compartments) are the most commonly affected sites. Causes include tibial fractures (most common cause of lower leg compartment syndrome), forearm fractures (Colles, Smith, radial neck), crushcrush injuries, vascular injuries, burns, limb compression, and reperfusion injury after revascularisation. The earliest and most important symptom is severe, unremitting pain out of proportion to the injury, worse with passive stretching. Diagnosis is clinical, with compartment pressure measurement confirming the diagnosis.

\medterm{Congenital Scoliosis}-- A lateral curvature of the spine caused by congenital vertebral anomalies that arise during embryogenesis (failure of formation or segmentation). Failure of formation: hemivertebra (wedge-shaped vertebra causing progressive curvature). Failure of segmentation: unilateral unsegmented bar (block of bone on one side). Mixed: combination of hemivertebra and unsegmented bar (most progressive). Presents with a progressive spinal curvature noted in infancy or early childhood. Associated with intraspinal anomalies (tethered cord, syringomyelia, diastematomyelia in 20-30\%), renal anomalies (10-20\%), and cardiac defects (10-15\%). Diagnosis: spine X-ray and MRI (to evaluate intraspinal anatomy) and renal ultrasound. Treatment: observation for mild curves; bracing is generally ineffective for congenital scoliosis. Surgical correction: in situ fusion, hemivertebra excision, or growth-friendly techniques (growing rods, VEPTR) for young children with progressive curves.

\medterm{Cubital Tunnel Syndrome}-- Cubital tunnel syndrome: ulnar nerve compression at elbow (cubital tunnel: Osborne's ligament, FCU aponeurosis, medial epicondyle, olecranon). Clinical: paresthesias ulnar 1.5 digits, nocturnal, Tinel at elbow, Froment sign (IP flexion of thumb with key pinch), Wartenberg sign (abduction of small finger), thenar/hypothenar wasting late. McGowan classification: I (intermittent), II (constant sensory, no motor), III (motor weakness). NCS/EMG: slowing across elbow. Treatment: I/II: nonoperative (ninight orthosis (elbow splint at 45-60 degrees of flexion), activity modification, ergonomics; Grade III or failed nonop: cubital tunnel release (in situ decompression, medial epicondylectomy, or transposition).

\medterm{Degenerative Spine Disease}-- Lumbar spinal stenosis: central (neurogenic claudication: bilateral leg pain with walking, relieved by sitting/flexion), lateral recess, foraminal. Imaging: MRI (central canal AP diameter <10mm, cross-sectional area <70mm2). Treatment: nonoperative (PT, epidural steroid, prostaglandin); operative: laminectomy (open, minimally invasive, tubular), decompression + fusion if instability (spondylolisthesis, scoliosis). Herniated nucleus pulposus: posterolateral (radiculopathy), central (cauda equina:cauda equina syndrome: surgical emergency; treated with urgent decompression.

\medterm{De Quervain Tenosynovitis}-- De Quervain tenosynovitis: first dorsal compartment (APL, EPB) stenosis. Clinical: radial wrist pain, Finkelstein test (thumb in palm, ulnar deviation), swelling over radial styloid. Treatment: splint (thumb spica), NSAIDs, corticosteroid injection (80-90\% cure); surgical: first dorsal compartment release (identify radial sensory nerve branches, watch for subcompartments).

\medterm{Developmental Dysplasia of the Hip}-- A spectrum of hip abnormalities in infants and children, ranging from mild acetabular dysplasia to complete hip dislocation. The femoral head is partially or completely displaced from the acetabulum due to capsular laxity and acetabular underdevelopment. Risk factors: female sex (4:1), breech presentation, family history, first-born, and oligohydramnios. Screening: clinical examination at every well-child visit (Ortolani test for reduction, Barlow test for dislocation, Galeazzi sign, asymmetric thigh/gluteal folds). Imaging: hip ultrasound for infants <4-6 months (dynamic and static views); X-ray for older infants and children (Hilgenreiner line, Perkin line, acetabular index). Treatment: Pavlik harness for infants <6 months (success rate 85-95\%); closed reduction and spica casting for harness failure or older infants; open reduction and pelvic/femoral osteotomy for children >18 months or failed closed reduction. Complications: avascular necrosis of the femoral head, redislocation, and early osteoarthritis.

\medterm{Diskectomy}-- The surgical removal of herniated intervertebral disk material that compresses nerve roots or the spinal cord, causing radiculopathy or myelopathy. Indicated in severe disk herniation causing unremitting pain, progressive motor weakness, or cauda equina syndrome.

\medterm{Dislocation}-- Complete displacement of the articular surfaces of a joint, resulting in loss of normal bony alignment and often associated with ligamentous, capsular, and neurovascular injury. The most commonly dislocated joints are the glenohumeral shoulder joint (accounting for 50 percent of all dislocations, anterior in 95 percent of cases), followed by the interphalangeal joints of the fingers, patellofemoral joint, elbow, and hip. The mechanism typically involves traumatic force applied to the joint beyonbeyond the normal range of motion, usually with significant force, causing capsular and ligamentous tearing, and often associated fractures, nerve injuries (axillary nerve in anterior shoulder dislocation, radial nerve in posterior shoulder dislocation), vascular injuries, and soft tissue injuries. Treatment involves urgent reduction followed by immobilization and rehabilitation.

\medterm{Distal Radius Fractures}-- The most common upper extremity fracture. Classification: AO/OTA: A (extra-articular), B (partial articular), C (complete articular); Frykman. Colles fracture: dorsal angulation/displacement (fall on outstretched hand). Smith fracture: volar angulation (fall on flexed wrist). Barton fracture: intra-articular fracture-dislocation (dorsal or volar). Chauffeur fracture: radial styloid. Treatment: nondisplaced/extra-articular: closed reduction + cast/splint; displaced intdisplaced intra-articular: ORIF (volar locking plate, dorsal plate, external fixation, fragment-specific fixation). Post-op: splint, early motion (for ORIF). Complications: malunion, nonunion (rare), carpal tunnel syndrome, post-traumatic arthritis, EPL rupture (3-5 percent), and complex regional pain syndrome (CRPS type I).

\medterm{Dupuytren contracture}-- A benign, progressive fibroproliferative disorder of the palmar fascia causing nodules, cords, and contractures of the fingers, most commonly affecting the ring and small fingers. The condition is caused by proliferation of myofibroblasts and deposition of type III collagen in the palmar aponeurosis, leading to formation of pathologic cords that contract and flex the metacarpophalangeal (MCP) and proximal interphalangeal (PIP) joints. Prevalence is highest in individuals of Northern European desNorthern European descent (up to 30 percent in some Scandinavian populations), with men affected 3 to 4 times more than women. The condition is associated with diabetes mellitus, epilepsy, alcoholism, and manual labour. The diagnosis is clinical.

\medterm{Ectrodactyly}-- A congenital limb malformation characterized by a deep median cleft of the hand or foot with missing central digits, resulting in a claw-like or lobster-claw appearance. Also known as split hand/foot malformation (SHFM). Most cases are autosomal dominant with incomplete penetrance, linked to mutations in TP63, DLX5, or WNT10B genes. May be isolated or syndromic (EEC syndrome — ectrodactyly, ectodermal dysplasia, cleft lip/palate; Split hand/foot malformation with long bone deficiency). Physical exam shows variable absence of the central rays (second, third, and/or fourth digits), ranging from a mild cleft to complete absence of central structures. Treatment: surgical reconstruction to improve pinch and grip function; prosthetics for severe cases. Genetic counseling is recommended.

\medterm{Ewing Sarcoma}-- A highly malignant small round blue cell tumor of bone and soft tissue, representing the second most common primary bone malignancy in children and adolescents (peak age 10-20 years), with a strong predilection for males and Caucasian populations. It arises from the EWSR1-FLI1 fusion gene (t(11;22)(q24;q12)) present in approximately 85 percent of cases, with other EWSR1 translocations (EWSR1-ERG, EWSR1-WT1) in the remainder. Unlike osteosarcoma, Ewing sarcoma most commonly involinvolves the diaphysis of long bones (femur, tibia, humerus) and the pelvis, with a predilection for the femur (25 percent), pelvis (20 percent), and tibia (15 percent). The typical presentation is painful swelling, often with fever, warmth, and erythema. Diagnosis is by biopsy and molecular genetics showing the EWSR1-FLI1 fusion transcript.

\medterm{Femoral Shaft Fractures}-- Femoral shaft fractures: high energy (MVC, fall from height) or low energy (elderly, pathologic). Winquist-Hansen classification of comminution: I (none), II (butterfly <50\% width), III (butterfly >50\%), IV (segmental). Treatment: antegrade intramedullary nailing (gold standard, reamed, statically locked), retrograde nailing (pregnant, bilateral, ipsilateral tibia fracture, morbid obesity), plate fixation (distal metaphyseal, nonunion). Early fixation (<24h) reduces ARDS, MODS, mortality (EPIC).(EPIC, early versus delayed stabilisation of the femur in multiply injured patients, with damage control orthopedics involving initial temporary external fixation followed by delayed definitive nailing at 48 to 72 hours, being the current standard of care for the polytrauma patient).

\medterm{Fracture}-- A break in the continuity of bone, classified by location, pattern (transverse, oblique, spiral, comminuted), displacement, and whether the fracture communicates with the external environment (open vs. closed). Fracture healing occurs through inflammatory, reparative, and remodeling phases. Management: reduction, immobilization (cast, splint, or surgical fixation with plates/screws/intramedullary nails), and rehabilitation. Open fractures require urgent debridement and antibiotics. Complications: nonunion, malunion, infection, compartment syndrome, and avascular necrosis.

\medterm{Ganglion Cysts}-- Ganglion cysts: mucin-filled cysts from joint capsule/tendon sheath. Dorsal wrist (60-70\%, scapholunate interval), volar wrist (radial artery), flexor tendon sheath (volar retinaculum, "seed ganglion"), DIP (mucous cyst, Heberden node association). Clinical: fluctuant mass, transilluminates, may cause nerve compression (volar: median/ulnar). Treatment: observation (spontaneous resolution 30-50\%), aspiration (high recurrence), surgical excision (capsule stalk removal, recurrence 5-10\%). Mucous cyMucous cyst (ganglion at DIP joint): associated with Heberden node, risk of nail deformity. Treatment: decompression, excision, pinning of DIP joint.

\medterm{Golfer Elbow}-- Medial epicondylitis, a tendinopathy of the common flexor/pronator tendon at the medial epicondyle. Less common than tennis elbow. Caused by repetitive wrist flexion and forearm pronation. Presents with medial elbow pain worsened by wrist flexion and gripping. Tenderness over the medial epicondyle. Positive Golfer elbow test. Risk of ulnar nerve compression. Diagnosis: clinical. Treatment: rest, ice, activity modification, NSAIDs, physical therapy (eccentric exercises), and corticosteroid injection (avoid ulnar nerve). Refractory cases: surgical release.

\medterm{Hallux Valgus}-- Lateral deviation of the great toe (hallux) with medial prominence of the first metatarsal head (bunion). More common in women, associated with narrow-toed shoes, high heels, and genetic predisposition. Presents with pain over the medial eminence, difficulty with shoe wear, and progressive deformity. Second toe can become hammer toe from crowding. Diagnosis: weight-bearing foot X-ray (hallux valgus angle >15°, intermetatarsal angle >9°). Treatment: wide-toed shoes, orthotics, toe spacers, and NSAIDs. Surgery (osteotomy, arthrodesis, or arthroplasty) for severe deformity or failed conservative treatment.

\medterm{Hallux Valgus (Bunion)}-- A progressive deformity of the first metatarsophalangeal (MTP) joint, characterized by lateral deviation of the great toe (hallux valgus) and medial deviation of the first metatarsal head (metatarsus primus varus), resulting in a bony prominence on the medial side of the forefoot (the bunion). Epidemiology: 3 times more common in women, prevalence increases with age (up to 35\% in adults >65). Risk factors: narrow-toed shoes with high heels (most significant modifiable risk factor), family history, heredity, pes planus (flatfoot), hypermobile first ray, ligamentous laxity, and inflammatory arthritis (rheumatoid arthritis, gout). Pathophysiology: the proximal phalanx rotates and drifts laterally, the first metatarsal head drifts medially, the sesamoids sublux laterally, and the MTP joint capsule stretches medially. Symptoms: pain over the medial eminence (bursitis), difficulty fitting shoes, cosmetic concern, and progression of deformity. May be associated with hammer toe of the second toe and transfer metatarsalgia. Treatment: (1) Conservative: footwear modification (wide-toed shoes, soft uppers, bunion pads/orthoses/splints; NSAIDs for acute pain). (2) Surgical: over 150 described techniques - distal chevron osteotomy (mild-moderate), Scarf osteotomy (moderate-severe), Lapidus arthrodesis (severe with first TMT instability), Akin osteotomy (proximal phalanx). Post-op: 6 weeks non-weight-bearing in a postoperative shoe or boot, gradual return to normal footwear.

\medterm{Hammer Toe}-- A deformity of the second, third, or fourth toe characterized by flexion at the PIP joint and extension at the MTP and DIP joints. Caused by muscle imbalance (intrinsic foot muscle weakness) and narrow-toed shoes. Presents with pain over the dorsal PIP joint (corn/callus) and the tip of the toe. Flexible hammer toe: correctable passively. Rigid: fixed deformity. Diagnosis: clinical. Treatment: conservative: wide toe box shoes, toe crests/pads, and stretching exercises. Surgical: flexor tenotomy (flexible) or PIP arthroplasty/arthrodesis (rigid).

\medterm{Hamstring}-- One of three muscles forming the posterior compartment of the thigh: semitendinosus, semimembranosus, and biceps femoris (long and short heads). Common origin: ischial tuberosity (except short head of biceps femoris which originates from the linea aspera of the femur). Insertion: semitendinosus and semimembranosus to the medial tibial condyle (pes anserinus area); biceps femoris to the fibular head. Actions: hip extension, knee flexion, and (with knee flexed) tibial rotation (medial by semitendinosus/semimembranosus, lateral by biceps femoris). Innervation: tibial branch of sciatic nerve (L5-S2) for all except short head of biceps femoris (common peroneal nerve, L5-S2). Hamstring strain is the most common muscle injury in athletes (especially sprinters, football, rugby), typically involving the proximal musculotendinous junction of biceps femoris. Grades: grade 1 (mild, microscopic tears), grade 2 (partial tear), grade 3 (complete rupture). Risk factors: eccentric overload, fatigue, poor flexibility, previous hamstring injury, muscle imbalance (quadriceps dominance). Management: RICE (rest, ice, compression, elevation) in acute phase, followed by progressive rehabilitation (eccentric strengthening, Nordic hamstring curl for prevention, core stability, sport-specific training). Return to sport when full range of motion, strength >90\% of unaffected side, and sport-specific movements are pain-free.

\medterm{Herniated Disc}-- Protrusion of the nucleus pulposus through a tear in the annulus fibrosus of the intervertebral disc, causing compression or inflammation of adjacent nerve roots. Most common at L4-L5 and L5-S1 (lumbar) and C5-C6, C6-C7 (cervical). Lumbar: sciatica (radiating leg pain), paresthesias, sensory loss, and motor weakness (foot drop). Cervical: neck pain radiating to arm with sensory/motor deficits. Cauda equina syndrome: surgical emergency with saddle anesthesia, bowel/bladder dysfunction, and bilateral leg weakness. Diagnosis: MRI. Treatment: conservative first (NSAIDs, PT, activity modification, epidural steroid injections). Surgical: microdiscectomy for refractory pain, progressive weakness, or cauda equina syndrome.

\medterm{Herniated disk}-- Also known as intervertebral disc herniation, displacement of nucleus pulposus material through a tear in the annulus fibrosus, most commonly occurring in the lumbar spine (L4-L5 and L5-S1 account for 95 percent of lumbar herniations) and cervical spine (C5-C6 and C6-C7). The pathophysiology involves age-related degenerative changes in the intervertebral disc: desiccation of the nucleus pulposus, fissuring of the annulus fibrosus, and eventual extrusion of nuclear material into the spinal canal or neural foramen, where the herniated disc material can compress the spinal cord (central herniation) or the nerve roots (posterolateral herniation, the most common), causing the characteristic radicular pain, sensory and motor deficits, and reflex changes in a dermatomal and myotomal distribution.

\medterm{Hip Disorders}-- The hip joints are ball-and-socket joints connecting the femur to the pelvis, providing stability and a wide range of motion for walking, running, and weight-bearing activities. Common hip conditions include osteoarthritis, fractures, bursitis, and labral tears. Management involves physical therapy, pain management, and joint replacement surgery for advanced degenerative disease.

\medterm{Hip Fractures}-- Common in elderly (osteoporosis, falls), classified by location: femoral neck (intracapsular), intertrochanteric (extracapsular), subtrochanteric. Garden classification for femoral neck: I (incomplete/impacted), II (complete nondisplaced), III (complete partially displaced), IV (complete displaced). Treatment: Garden I/II: cannulated screws; Garden III/IV in elderly: hemiarthroplasty (bipolar preferred) or total hip arthroplasty (if active, pre-existing arthritis); young patienpatients (<65, good bone): ORIF (cannulated screws or sliding hip screw, blade plate, or cephalomedullary nail depending on fracture pattern). Subtrochanteric fractures: high nonunion rate, IM nail (long, statically locked, reamed) versus plate fixation.

\medterm{Hip Replacement}-- A surgical procedure in which the diseased hip joint is replaced with a prosthetic implant (arthroplasty). Total hip replacement (THR) replaces both the femoral head and the acetabulum. Indications: osteoarthritis (most common, 60\%), rheumatoid arthritis, avascular necrosis (AVN) of the femoral head, femoral neck fracture (especially displaced in elderly), hip dysplasia, ankylosing spondylitis, Paget disease, and failed previous hip surgery. Components: acetabular component (metal shell with polyethylene or ceramic liner), femoral component (metal stem with modular head--cobalt-chrome or ceramic). Fixation: cemented (using polymethylmethacrylate bone cement, preferred in elderly with osteoporotic bone) or cementless (press-fit, with porous coating for bone ingrowth, preferred in younger patients with good bone quality). Surgical approaches: posterior (most common, risk of dislocation 2-5\% in first 3 months, repairs capsule/external rotators), lateral (Hardinge, gluteus medius split, lower dislocation risk), anterolateral (Watson-Jones), and direct anterior (muscle-sparing, faster recovery, less postoperative pain). Outcomes: excellent pain relief and functional improvement in 90-95\% at 10 years, 80-85\% implant survival at 20 years. Complications: dislocation (2-5\%, highest in first 3 months), periprosthetic fracture, infection (1-2\%, devastating, requires two-stage revision: removal of implant + antibiotic cement spacer + 6 weeks IV antibiotics + reimplantation), aseptic loosening (most common cause of late failure, due to wear debris-induced osteolysis, polyethylene wear particles stimulate macrophage activation and RANKL-mediated osteoclastic bone resorption), leg length discrepancy (subjective or objective), DVT/PE (routine thromboprophylaxis), nerve injury (sciatic, femoral).

\medterm{Joint Diseases}-- A broad category of pathological conditions affecting the structure and function of joints. Classification by etiology: (1) Degenerative: osteoarthritis (most common, non-inflammatory, affects load-bearing joints--knee, hip, spine). (2) Inflammatory: rheumatoid arthritis (autoimmune, symmetric small-joint synovitis), psoriatic arthritis, ankylosing spondylitis, reactive arthritis, gout (crystal arthropathy, uric acid, podagra), pseudogout (CPPD). (3) Infectious: septic arthritis (bacterial, urgent, joint destruction within 24 hours), Lyme arthritis, tuberculosis arthritis, viral arthritis (parvovirus, rubella, hepatitis B/C). (4) Metabolic: gout, pseudogout, haemochromatosis (chondrocalcinosis, calcium pyrophosphate), ochronosis (alkaptonuria, homogentisic acid accumulation, pigmented cartilage, spondylosis). (5) Trauma: intra-articular fracture, dislocation, ligament tear, meniscal tear, haemarthrosis. (6) Neoplastic: synovial sarcoma, pigmented villonodular synovitis (benign, recurrent haemarthrosis). (7) Childhood: juvenile idiopathic arthritis (JIA). Diagnosis: history (mechanical vs inflammatory pain, stiffness, swelling), examination (warmth, effusion, crepitus, range of motion, instability), imaging (X-ray, ultrasound, MRI), synovial fluid analysis (cell count, crystals, Gram stain, culture). Diagnosis: medical history, physical examination, imaging (X-ray, ultrasound, MRI), and synovial fluid analysis. Treatment directed at the underlying cause.

\medterm{Joint Pain}-- Discomfort arising from articular structures including cartilage, synovium, ligaments, and periarticular soft tissues due to inflammation, degeneration, injury, or infection. Common causes: osteoarthritis (degenerative, most prevalent), rheumatoid arthritis (autoimmune inflammatory), gout (crystal arthropathy), septic arthritis (bacterial infection requiring urgent intervention), and trauma. Clinical features include pain worsened by movement, stiffness, swelling, warmth, and decreased range of motion. Diagnosis requires history, physical exam, imaging (X-ray, MRI), and laboratory studies (inflammatory markers, autoantibodies, uric acid, synovial fluid analysis). Treatment targets the underlying cause: NSAIDs or corticosteroids for inflammation, disease-modifying antirheumatic drugs (DMARDs) for inflammatory arthritis, analgesics, physical therapy, and joint replacement for advanced disease.

\medterm{Joint Replacement (Arthroplasty)}-- A surgical procedure in which a damaged or arthritic joint surface is removed and replaced with a prosthetic implant, most commonly performed for the hip (total hip arthroplasty) and knee (total knee arthroplasty). Indications include end-stage osteoarthritis, rheumatoid arthritis, avascular necrosis, and severe joint trauma causing debilitating pain and functional limitation refractory to conservative therapy. The procedure involves resurfacing the articular ends of bone with metal, polyethylene, or ceramic components fixated with cement or press-fit technique. Postoperative management includes early mobilization, physical therapy, thromboprophylaxis, and activity restrictions to prevent dislocation. Complications: infection, loosening, fracture, nerve injury, venous thromboembolism, and prosthetic wear requiring revision.

\medterm{Joints}-- The sites where two or more bones meet, providing structural support and enabling movement. Classification by structure and function: (1) Fibrous joints (synarthroses, immovable): sutures of the skull, gomphoses (tooth in socket), syndesmoses (tibia-fibula). (2) Cartilaginous joints (amphiarthroses, slightly movable): synchondroses (costochondral, growth plates), symphyses (pubic symphysis, intervertebral discs). (3) Synovial joints (diarthroses, freely movable): characterized by a joint cavity lined by synovial membrane containing synovial fluid, articular cartilage (hyaline), and a fibrous joint capsule. Further classified by shape: hinge (elbow, knee), ball-and-socket (hip, shoulder), pivot (atlantoaxial, proximal radioulnar), condylar (temporomandibular, knee--also hinge), saddle (first carpometacarpal--thumb), gliding (intercarpal, intertarsal), and ellipsoid (wrist). Synovial fluid: plasma dialysate with hyaluronic acid (lubrication, shock absorption, nutrition for avascular articular cartilage). Joint health is maintained by low-impact exercise, weight-bearing activity, and avoidance of excessive repetitive loading. Joint pathology: osteoarthritis, rheumatoid arthritis, gout, septic arthritis, trauma (dislocation, fracture), and haemarthrosis.

\medterm{Kienbock Disease}-- Kienbock disease (lunatomalacia): AVN of lunate. Lichtman classification: I (normal XR, MRI edema), II (lunate sclerosis), IIIA (lunate collapse, no scaphoid rotation), IIIB (lunate collapse, fixed scaphoid rotation, DISI), IV (perilunate arthritis). Etiology: negative ulnar variance, trauma, vascular variation. Treatment: I/II: radial shortening (negative ulnar variance), capitate shortening, vascularized bone graft (4,5-extensor compartments artery); IIIA: radial shortening, capitolunate fusiocapitolunate fusion, capitate-lunate-triquetrum fusion, lunate core decompression; IIIB-IV: salvage (proximal row carpectomy, scaphoidectomy plus 4-corner fusion, total wrist arthrodesis).

\medterm{Klippel-Feil Syndrome}-- A congenital condition characterized by fusion of two or more cervical vertebrae, resulting from failure of normal segmentation of the cervical somites during weeks 3-8 of gestation. Classic triad: short neck, low posterior hairline, and limited cervical range of motion (present in <50\% of cases). Patients are at risk for atlantoaxial instability, cervical stenosis, and the CRUS sequence (Cervical fusion, Renal anomalies, Uterine anomalies, and Sensorineural hearing loss). Associated anomalies: Sprengel deformity (30\%), scoliosis (60\%), renal agenesis/ectopia, congenital heart disease (VSD, ASD), hearing loss, and synkinesia (mirror movements). Diagnosis: cervical spine X-ray or CT shows fused vertebral bodies and posterior elements; MRI for spinal cord and brainstem assessment; renal ultrasound and audiometry for associated anomalies. Treatment: activity restriction to prevent cervical injury; surgical stabilization for instability or neurologic compromise.

\medterm{Kneecap}-- See Patella (Kneecap).

\medterm{Knee Disorders}-- The knee is a complex hinge-type synovial joint connecting the femur, tibia, and patella, stabilized by ligaments (ACL, PCL, MCL, LCL), menisci (medial and lateral), and surrounding muscles/tendons. Common knee disorders include osteoarthritis (degenerative cartilage loss most prevalent in older adults), anterior cruciate ligament (ACL) tears (sports-related pivot injury), meniscal tears (twisting injury with mechanical locking or clicking), patellofemoral pain syndrome (anterior knee pain in active young adults), and Baker's cyst (popliteal swelling from fluid extrusion). Diagnosis relies on history, physical examination (Lachman test, McMurray test, joint line tenderness), and imaging (weight-bearing X-rays, MRI for soft-tissue pathology). Treatment ranges from conservative (activity modification, physical therapy, bracing, NSAIDs) to surgical (arthroscopic meniscectomy/repair, ligament reconstruction, partial or total knee arthroplasty).

\medterm{Kyphosis}-- An excessive anterior-posterior curvature of the thoracic spine, measured by the Cobb angle on lateral spine x-ray (normal thoracic kyphosis is 20-40 degrees). Types: postural kyphosis (most common, flexible, correctable with hyperextension, in adolescents, benign, not progressive, no structural vertebral abnormalities), Scheuermann kyphosis (rigid, structural kyphosis in adolescents, caused by anterior wedging of greater than 5 degrees in three or more consecutive thoracic vertebrae, with endplendplate irregularities, Schmorl nodes, and decreased disc height), and congenital kyphosis (rare, due to failure of formation or segmentation of vertebral bodies). The diagnosis is made by standing lateral x-ray with the Cobb angle measured from the superior endplate of one vertebra to the inferior endplate of the lowest involved vertebra.

\medterm{Lateral Collateral Ligament Injury}-- Tear of the LCL, less common than MCL injury but more significant. Mechanism: varus stress to the knee (direct blow to medial knee). Presents with lateral knee pain, swelling, and varus laxity. Part of the posterolateral corner (PLC) — when injured, causes significant rotational instability. Grade III LCL injury is often associated with PCL injury and peroneal nerve injury. Diagnosis: varus stress test at 0° and 30°; MRI. Treatment: Grade I-II: bracing. Grade III or multiligament: surgical repair or reconstruction.

\medterm{Lordosis}-- The normal inward curvature of the cervical and lumbar spine. Excessive lumbar lordosis (hyperlordosis): increased pelvic tilt, commonly due to hip flexor tightness, weak abdominals, or pregnancy. Presents with low back pain and forward protrusion of the abdomen and buttocks. Causes: postural, congenital (achondroplasia), neuromuscular (cerebral palsy), and compensatory (kyphosis). Diagnosis: clinical, standing lateral X-ray. Treatment: postural exercises, core strengthening, hamstring/hip flexor stretching. Surgery rarely needed.

\medterm{Lumbago}-- A non-specific term for low back pain, typically referring to acute or chronic pain in the lumbar region (between the lower ribs and the gluteal folds). Lumbago is a symptom, not a diagnosis. Causes: mechanical low back pain (most common, 90\%, includes lumbar strain/sprain, discogenic pain, facet joint pain, sacroiliac joint dysfunction), disc herniation with radiculopathy (sciatica, 5\%), spinal stenosis, spondylolisthesis, vertebral compression fracture (osteoporotic or traumatic), infection (discitis, osteomyelitis, epidural abscess), malignancy (bone metastases, multiple myeloma, primary bone tumor), inflammatory (ankylosing spondylitis, psoriatic arthritis), and visceral referred pain (pancreatitis, nephrolithiasis, AAA, endometriosis). Red flags (cauda equina syndrome, infection, malignancy, fracture): saddle anesthesia, bowel/bladder incontinence, bilateral sciatica, progressive neurological deficit, fever, night pain, weight loss, age >50, history of cancer, IV drug use, immunosuppression. Yellow flags (psychosocial risk factors for chronicity): fear-avoidance beliefs, catastrophising, depression, job dissatisfaction, pending litigation. Most acute lumbago resolves within 4-6 weeks. Management: conservative--analgesics (paracetamol, NSAIDs--ibuprofen, naproxen, diclofenac), muscle relaxants (short-term, diazepam, cyclobenzaprine), heat/ice, activity modification (avoid bed rest, stay active), physiotherapy (core stability, McKenzie exercises, manual therapy). For radicular pain: epidural corticosteroid injections. Surgery: discectomy for refractory sciatica, decompression for spinal stenosis, fusion for instability.

\medterm{Medial Collateral Ligament Injury}-- Tear of the MCL, the most common knee ligament injury. Mechanism: valgus stress to the knee (direct blow to lateral knee). Presents with medial knee pain, swelling, and valgus laxity. Grading: I (pain, no laxity), II (partial tear, some laxity), III (complete tear, significant laxity). Diagnosis: valgus stress test at 0° and 30° of flexion; MRI. Treatment: almost always conservative — functional bracing with early motion. Grade III injuries with associated ACL/PCL/PLC injuries may require surgical repair.

\medterm{Meniscal Injuries}-- Meniscal tears: traumatic (vertical longitudinal, bucket handle, radial, horizontal cleavage, flap) vs degenerative (complex, horizontal). Clinical: joint line tenderness, McMurray, Apley, Thessaly test. MRI: sensitivity 90\%+. Treatment: repairable (vertical longitudinal, peripheral 100 3mm, red-red/red-white zone, young): inside-out, all-inside, outside-in; irreparable (degenerative, white-white, complex): partial meniscectomy (preserve rim). Root tears: radial tears at posterior horn attachmenattachment site (medial meniscus). Treatment: repair (transtibial pull-out, suture anchor) vs meniscectomy.

\medterm{Meniscal Tear}-- Tear of the C-shaped fibrocartilage (medial or lateral meniscus) that provides shock absorption and stability in the knee. Medial meniscus tears more common (attached to capsule). Mechanism: twisting injury with weight-bearing. Presents with joint line tenderness, pain with deep squatting, catching, locking, and giving way. Positive McMurray test. Diagnosis: MRI. Treatment: conservative for small, peripheral (red zone) tears; arthroscopic partial meniscectomy for large, complex tears that cannot be repaired; meniscal repair for peripheral tears in young patients.

\medterm{Muscle Pain (Myalgia)}-- Myalgia—muscle pain—is a common symptom arising from skeletal muscle injury, inflammation, overuse, infection, or systemic disease. Acute myalgia is often due to trauma (muscle strain/tear, contusion), eccentric exercise (delayed-onset muscle soreness peaking 24--72 hours post-exertion), or viral infections (influenza, COVID-19, enterovirus). Chronic diffuse myalgia characterizes fibromyalgia (widespread pain with tender points, fatigue, sleep disturbance, cognitive fog—affecting 2--4\% of population, 80\% female) and polymyalgia rheumatica (morning stiffness and aching in shoulder/pelvic girdles in adults >50, ESR/CRP markedly elevated). Diagnosis: CK, ESR, CRP, TSH, vitamin D, autoimmune serologies, EMG, and muscle biopsy when indicated. Treatment: symptomatic (acetaminophen, NSAIDs, heat/ice, massage), treatment of underlying cause (antibiotics for infectious myositis, corticosteroids for polymyalgia rheumatica/idiopathic inflammatory myopathies, exercise and tricyclic antidepressants/gabapentinoids for fibromyalgia).

\medterm{Nonunion and Malunion}-- Nonunion is failure of a fractured bone to heal after an adequate period (generally 6-9 months for long bones) without clinical or radiographic evidence of progression toward union for 3 months. Classification: hypertrophic (viable ends, inadequate stability), atrophic (devitalized ends, poor biology), oligotrophic (intermediate). Risk factors: open fractures, severe comminution, inadequate fixation, infection, smoking, NSAIDs, steroids, diabetes, vascular disease, radiation. Diagnosis: persistepersistent pain, implant failure. Treatment: hypertrophic (increase stability with compression plating, exchange nailing, dynamization); atrophic (improve biology with bone graft, vascularized bone graft).

\medterm{Open Fractures}-- Fractures with communication between the fracture site and the external environment, classified by Gustilo-Anderson: Grade I (<1 cm wound, clean), Grade II (1-10 cm wound, moderate soft tissue injury), Grade IIIA (>10 cm wound, adequate soft tissue coverage), Grade IIIB (>10 cm wound, extensive soft tissue loss requiring flap coverage), Grade IIIC (any open fracture with arterial injury requiring repair). Management: ATLS protocol, antibiotics within 1 hour (Grade I: cefazolincefazolin 1-2g IV q8h; Grade II/III: cefazolin plus gentamicin; Grade III: penicillin for farm injuries, tetanus prophylaxis, thorough irrigation and debridement within 6 hours, fracture stabilisation (external fixation for severe contamination), wound closure by 7 days.

\medterm{Orthopedic Biologics}-- Bone graft substitutes: autograft (iliac crest gold standard, osteogenic/osteoinductive/osteoconductive, donor site morbidity), allograft (osteoconductive, some osteoinductive if DBM, disease transmission risk), DBM (demineralized bone matrix, osteoinductive), synthetic (calcium phosphate, calcium sulfate, bioactive glass - osteoconductive), BMPs (rhBMP-2, rhBMP-7, potent osteoinductive, swelling/ectopic bone cost), stem cells (BMAC, MSCs, adipose-derived), PRP (platelet-rich plasma, growth factgrowth factors, cytokines). Complications: infection, graft resorption, nonunion, heterotopic ossification. Applications: nonunion, delayed union, spinal fusion, osteochondral defects.

\medterm{Orthopedic Infections}-- Septic arthritis: see Infectious Disease. Osteomyelitis: see Infectious Disease. Prosthetic joint infection: see Periprosthetic Joint Infection. Diabetic foot infection: Wagner classification (0-5), University of Texas (grade + ischemia + infection). Treatment: offloading (total contact cast), debridement, vascular assessment (ABI, TBI, TcPO2), revascularization (endovascular, bypass), antibiotics (mild: oral; moderate/severe: IV), amputation (toe, ray, transmetatarsal, BKA, AKA). Charcot neuroaosteoarthropathy (Charcot foot): midfoot collapse (rocker bottom), neuropathic, treated with Charcot reconstruction (superconstruct fixation, intramedullary beam, exostectomy, TCC, and AFO).

\medterm{orthopedics}-- The branch of surgery concerned with the musculoskeletal system, including bones, joints, ligaments, tendons, muscles, and nerves. Orthopaedic surgeons diagnose and treat conditions such as fractures and dislocations, arthritis (osteoarthritis, inflammatory arthritis), spinal disorders (disc herniation, spinal stenosis, scoliosis), sports injuries (ACL tear, rotator cuff tear, meniscal injury), paediatric orthopaedic conditions (developmental dysplasia of the hip, clubfoot, Perthes disease), bone tumors, and infections (osteomyelitis, septic arthritis). Treatment modalities: conservative (cast/splint, physiotherapy, medications), minimally invasive (arthroscopy, percutaneous fixation) and open surgery (joint replacement—hip, knee, shoulder; spinal fusion; osteotomy; limb reconstruction). Subspecialties include trauma, arthroplasty, hand surgery, foot and ankle, spine, paediatric orthopedics, sports medicine, and orthopaedic oncology.

\medterm{Osgood-Schlatter Disease}-- A traction apophysitis of the tibial tuberosity at the insertion of the patellar tendon, caused by repetitive microtrauma during periods of rapid growth. It is a common cause of anterior knee pain in adolescents aged 10--15 years, more common in boys (ratio 3:1) and in active/sporty children (running, jumping sports—basketball, football, volleyball). Presentation: gradual onset of unilateral (30\% bilateral) anterior knee pain just below the kneecap, worsened by activity (running, jumping, kneeling, climbing stairs), improved by rest. Physical exam: tenderness and swelling over the tibial tuberosity, pain on resisted knee extension or squatting. X-ray may show irregularity or fragmentation of the tibial tuberosity apophysis, but diagnosis is clinical. Management: activity modification (reduce high-impact sports, not complete rest), ice after activity, NSAIDs, quadriceps and hamstring stretching/strengthening, patellar tendon strap, and reassurance (self-limiting condition that resolves with skeletal maturity—apophyseal closure at age 14--18). Surgery (ossicle excision, tibial tuberosity reduction) is rarely needed for persistent symptoms.

\medterm{Osteochondritis Dissecans}-- A condition characterized by separation of a fragment of articular cartilage and subchondral bone from the joint surface, most commonly in the knee (lateral aspect of the medial femoral condyle). Most common in adolescents and young adults. Etiology: repetitive microtrauma, ischemia, or genetic predisposition. Presents with vague joint pain, swelling, catching, locking, and giving way. Diagnosis: X-ray (lucent crater with loose fragment), MRI (fragment stability and viability). Treatment: stable lesions in skeletally immature: activity modification, immobilization, or drilling. Unstable or loose bodies: arthroscopic fixation or fragment excision.

\medterm{Osteochondroma}-- The most common benign bone tumor, accounting for approximately 20-35 percent of all benign bone tumors. It is characterized by a bony projection with a cartilage cap arising from the surface of a bone, continuous with the underlying cortex and medullary cavity. Solitary osteochondromas typically present in the second decade of life (ages 10-20 years) as a painless, firm, immobile mass near the knee (distal femur or proximal tibia, most common), proximal humerus, oror the proximal tibia or humerus. The characteristic radiographic appearance is a pedunculated or sessile bony excrescence arising from the metaphysis, pointing away from the adjacent joint, with a cartilage cap continuous with the host bone. Malignant transformation to secondary chondrosarcoma occurs in approximately 1 to 2 percent of solitary osteochondromas but in 2 to 5 percent of patients with multiple hereditary exostosis.

\medterm{Osteoid Osteoma}-- A benign bone-forming tumor that accounts for approximately 10-12 percent of all benign bone tumors, predominantly affecting adolescents and young adults (ages 5-25 years, male-to-female ratio 2:1). It is characterized by a small (< 1.5-2 cm) radiolucent nidus surrounded by reactive sclerotic bone. The nidus contains osteoblasts, osteoid, and woven bone with prominent vascular channels. Osteoid osteomas produce high levels of prostaglandin E2, which causes the characteristic scharacteristic severe, localised, deep, boring pain that is worse at night and dramatically relieved by NSAIDs, which inhibit COX-2 and reduce the production of prostaglandin E2, the key mediator of pain in this tumor.

\medterm{Osteopathy}-- A system of manual medicine and healthcare founded by Andrew Taylor Still (1828--1917) based on the principle that the body's structure (musculoskeletal system) and function are interrelated, and that the body has inherent self-regulatory and healing mechanisms. Osteopathic practitioners (osteopaths in the UK; Doctors of Osteopathic Medicine, DOs, in the US) use hands-on techniques including soft tissue massage, joint articulation, myofascial release, high-velocity low-amplitude thrust, muscle energy technique, and cranial osteopathy to diagnose and treat musculoskeletal pain, dysfunction, and associated conditions. In the UK, osteopathy is regulated by the General Osteopathic Council (GOsC) and is a statutory regulated profession; osteopaths complete a 4-year integrated Master's degree (MOst). Osteopathy is distinct from orthopaedic surgery but often works alongside it in managing back pain, neck pain, and sports injuries. Evidence supports osteopathic manipulative treatment for chronic low back pain, although the evidence base for cranial osteopathy is limited.

\medterm{Osteopenia}-- Bone mineral density (BMD) that is lower than normal but not low enough to be classified as osteoporosis. T-score between -1.0 and -2.5 on DXA scan. A risk factor for osteoporosis and fractures. Management: lifestyle modifications (weight-bearing exercise, adequate calcium [1200 mg/day] and vitamin D [800-1000 IU/day]), smoking cessation, fall prevention, and reassess BMD every 1-2 years. Pharmacotherapy considered if FRAX score indicates high fracture risk.

\medterm{Osteopetrosis}-- A rare inherited bone disease characterized by defective osteoclast function, leading to impaired bone resorption and abnormally dense, sclerotic bones. The bone is structurally abnormal—brittle despite increased density—and prone to fracture. Two main forms: (1) Autosomal recessive ('malignant'): presents in infancy with failure to thrive, anemia, thrombocytopenia (bone marrow failure from marrow space obliteration), hepatosplenomegaly (extramedullary haematopoiesis), cranial nerve palsies (optic atrophy, deafness, facial palsy due to nerve compression), and recurrent infections. Life expectancy is poor without treatment. (2) Autosomal dominant ('benign'): milder, presents in adolescence or adulthood with fractures, osteomyelitis of the mandible, and sometimes asymptomatic (incidental finding on X-ray). Diagnosis: X-ray shows generalised osteosclerosis ('sandwich vertebrae', 'bone-within-bone' appearance, Erlenmeyer flask deformity of femurs). Treatment: hematopoietic stem cell transplantation (HSCT) for the malignant form (restores osteoclast function, improves survival), supportive care for fractures and anemia, and interferon gamma-1b. Gene therapy is experimental.

\medterm{Osteoporosis}-- See geriatrics.

\medterm{Osteosarcoma}-- The most common primary malignant bone tumor in children and adolescents (excluding leukemia), with a peak incidence during the adolescent growth spurt (ages 10-20 years), and a second smaller peak in older adults (usually secondary to Paget disease or prior radiation). It most commonly arises in the metaphysis of long bones, particularly the distal femur (most common, 30-40 percent), proximal tibia, and proximal humerus. Histologically, it is characterized by malignant osteoid pproduction (malignant osteoid) by malignant tumor cells. Clinical presentation is with localised pain (the most common symptom, worse at night), swelling, a palpable mass, and pathological fracture in 10 to 15 percent of cases.

\medterm{Osteotomy}-- A surgical procedure involving cutting and repositioning of bone to correct deformity, realign joint surfaces, alter weight-bearing forces, or change limb length. Indications: (1) Lower limb—high tibial osteotomy (HTO) for medial compartment knee osteoarthritis in young, active patients (realigns the mechanical axis, delaying arthroplasty), distal femoral osteotomy for lateral compartment or valgus deformity, and femoral/tibial osteotomy for congenital deformities (blount disease, leg length discrepancy). (2) Hip—pelvic osteotomy (periacetabular osteotomy—Ganz procedure) for developmental dysplasia of the hip in young adults, and femoral varus/valgus osteotomy for Perthes disease or hip dysplasia. (3) Spine—spinal osteotomy (Smith-Petersen, pedicle subtraction osteotomy, vertebral column resection) for sagittal imbalance (kyphosis, flatback syndrome). (4) Foot—calcaneal osteotomy for hindfoot deformity, first metatarsal osteotomy for hallux valgus. Osteotomy is also used in orthognathic surgery (jaw osteotomy for malocclusion). Fixation: plates, screws, external fixators. Recovery involves non-weight-bearing or protected weight-bearing for 6--12 weeks post-operatively.

\medterm{Paget Disease of Bone}-- A chronic bone disorder characterized by excessive, disorganized bone remodeling with increased osteoclastic bone resorption followed by compensatory osteoblastic bone formation, producing structurally abnormal bone that is enlarged, weak, and prone to fracture. Often asymptomatic; can present with bone pain, deformity (bowing of tibia, skull enlargement), fractures, and arthritis. Diagnosis: elevated serum alkaline phosphatase with normal calcium, phosphate, and PTH; X-ray findings (lytic lesions, bone enlargement). Treatment: bisphosphonates (alendronate, risedronate, zoledronic acid).

\medterm{Paget's Disease}-- See Paget Disease of Bone.

\medterm{Patella (Kneecap)}-- The largest sesamoid bone in the human body, embedded within the quadriceps tendon anterior to the knee joint. The patella increases the lever arm of the quadriceps muscle (improving mechanical advantage for knee extension by ~30\%), protects the anterior knee joint, and distributes compressive forces across the femoral trochlea. Articular surface: the posterior patella has a V-shaped ridge that articulates with the femoral trochlear groove (medial and lateral facets). Patellofemoral joint: transmits up to 3--5 times body weight during stair climbing and 7--8 times during squatting. Patellar dislocation: typically lateral (vastus medialis obliquus disruption, trochlear dysplasia, patella alta, increased Q-angle), reduced by extending the knee. Patellar fracture: caused by direct trauma or forceful quadriceps contraction (transverse fracture most common); non-displaced: cast immobilisation; displaced/disrupted extensor mechanism: tension band wiring. Patellofemoral pain syndrome (anterior knee pain): common in young active individuals and runners, treated with quadriceps strengthening, activity modification, and patellar taping.

\medterm{Patellofemoral Instability}-- Patellar dislocation: lateral dislocation (95\%), often reduces spontaneously. Risk factors: trochlear dysplasia (Dejour classification), patella alta (Caton-Deschamps >1.2), increased TT-TG distance (>20mm), MPFL insufficiency, femoral anteversion, genu valgum. First-time: nonoperative (brace, PT, MPFL healing) if no loose body; recurrent: MPFL reconstruction (gold standard), tibial tubercle osteotomy (Fulkerson anteromedialization) for patella alta/TT-TG >20mm, trochleoplasty (deepening) for sesevere trochlear dysplasia (Dejour type B/D).

\medterm{Patellofemoral Pain Syndrome}-- Anterior knee pain related to the patellofemoral joint, common in adolescents and young adults (especially runners). Etiology: multifactorial (malalignment, quadriceps imbalance, overuse, tight lateral structures, weak VMO, and patellar instability). Presents with retropatellar or peripatellar pain with activities (squatting, running, stair climbing, prolonged sitting with knees bent — theater sign). Diagnosis: clinical; X-ray and MRI to rule out other pathology. Treatment: quadriceps and VMO strengthening, IT band and hamstring stretching, patellar taping/bracing, and activity modification. Most improve with non-operative management.

\medterm{PCL Injuries}-- Posterior cruciate ligament (PCL) tears: dashboard injury (posterior tibial force), hyperflexion. Clinical: posterior drawer (most sensitive), quadriceps active test, posterior sag sign. Grade I (1-5mm), II (6-10mm), III (>10mm). Isolated PCL: nonoperative (brace, rehab) for Grade I/II; surgical for Grade III or combined. PCL reconstruction: transtibial or tibial inlay (open posterior approach), single or double bundle. Combined with PLC/ACL/MCL for multi-ligament.

\medterm{Peripheral Nerve Injuries}-- Peripheral nerve injury classification (Sunderland): 1st degree (neuropraxia, conduction block, full recovery), 2nd degree (axonotmesis, Wallerian degeneration, endoneurium intact, recovery), 3rd degree (endoneurium disrupted, perineurium intact, partial recovery), 4th degree (perineurium disrupted, neuroma-in-continuity, no recovery without surgery), 5th degree (neurotmesis, complete transection). Management: acute sharp transection: primary epineurial repair (9-0/10-0 nylon); blunt/crush: waitwait 3-6 weeks, then repair; chronic (>6-12 months): nerve transfer, tendon transfer, free functional muscle transfer. Prognosis: better for distal, sharp, clean injuries; worse for proximal, blunt, contaminated.

\medterm{Phocomelia}-- A severe congenital limb deficiency where the proximal portion of a limb (humerus/femur, radius/tibia) is absent or severely hypoplastic, and the hand or foot attaches directly to the trunk. Most famously associated with thalidomide exposure during the first trimester of pregnancy (days 20-36 after conception). Can also occur in Roberts syndrome, Holt-Oram syndrome, and TAR syndrome. Complete phocomelia: hand/foot directly attached to the trunk. Proximal phocomelia: absent humerus/femur with forearm/leg present. Distal phocomelia: absent radius/ulna with humerus present. Treatment: prosthetics and adaptive devices to optimize function; reconstructive surgery for selected cases. Genetic counseling and prenatal diagnosis are important.

\medterm{Plantar Fasciitis}-- The most common cause of heel pain, caused by degeneration and microtears of the plantar fascia at its calcaneal origin. Presents with sharp, stabbing heel pain with the first steps in the morning (post-static dyskinesia), improving after a few minutes but worsening with prolonged standing or activity. Risk factors: tight Achilles/gastrocnemius, obesity, prolonged standing, and high-impact activities. Diagnosis: clinical. Treatment: calf and plantar fascia stretching (first-line), orthotics (arch support, heel cup), night splints, NSAIDs, and corticosteroid injection. Refractory: shockwave therapy, PRP, or fasciotomy.

\medterm{Polydactyly}-- A congenital anomaly characterized by supernumerary digits on the hands or feet. Postaxial polydactyly (ulnar/fibular side, most common): more frequent in African American and Native American populations, often isolated and autosomal dominant. Preaxial polydactyly (radial/tibial side): more common in Asian populations, may be associated with syndromes (Holt-Oram, Fanconi anemia, VACTERL, TAR syndrome). Central polydactyly: rare, often syndromic. Classification: Type A (well-formed, articulating digit) vs. Type B (rudimentary, pedunculated skin tag). Physical exam and X-ray to assess for bony connection, joint articulation, and functional impact. Treatment: surgical excision for aesthetic and functional concerns; Type B can be ligated in the nursery. Preaxial polydactyly may require more complex reconstruction.

\medterm{Posterior Cruciate Ligament Injury}-- Tear of the PCL, less common than ACL injury. Mechanism: dashboard injury (tibia striking dashboard), hyperflexion, or fall on a flexed knee with foot plantarflexed. Presents with posterior knee pain, mild swelling, and posterior sag of the tibia. Positive posterior drawer test (most sensitive), posterior sag sign (Godfrey test). Often associated with PLC injury (posterolateral corner: popliteus, popliteofibular ligament, LCL). Diagnosis: MRI. Treatment: Grade I-II: conservative with quadriceps strengthening. Grade III or combined injuries: PCL reconstruction.

\medterm{Proximal Humerus Fractures}-- Proximal humerus fractures: Neer classification (4 parts: articular surface, greater tuberosity, lesser tuberosity, shaft; displaced >1 cm or angulated >45 degrees = part). 1-part (nondisplaced): sling, early motion. 2-part: surgical neck (ORIF plate, IM nail) or greater tuberosity (if displaced >5mm or rotator cuff involvement). 3-part: ORIF (locking plate) or hemiarthroplasty (elderly, osteoporotic). 4-part: high risk of AVN; hemiarthroplasty or reverse total shoulder arthroplasty (RSA, preferpreferred, particularly for elderly patients with 4-part fractures). Post-op: sling, early passive ROM, active assist at 6 weeks, strengthening at 12 weeks. Complications: AVN (up to 30 percent for 4-part), malunion, nonunion, stiffness, hardware failure, infection.

\medterm{Radial Club Hand}-- A congenital longitudinal deficiency of the radial ray, ranging from thumb hypoplasia to complete absence of the radius. The hand is radially deviated due to the absent radial support, with the carpus displaced proximally and radially. Associated with Holt-Oram syndrome (TBX5 mutation, with VSD or ASD), TAR syndrome (thrombocytopenia-absent radius, bilateral radial aplasia with thumbs present — key distinguishing feature), Fanconi anemia (progressive pancytopenia, radial anomalies), and VACTERL association. Physical exam shows radial deviation of the hand, short forearm, thumb anomalies (absent, hypoplastic, or triphalangeal), and limited elbow extension. X-ray confirms radius deficiency and associated skeletal anomalies. Treatment: centralization or radialization of the wrist (surgical repositioning over the ulna), thumb reconstruction (pollicization), and splinting. Genetic testing for Fanconi anemia is essential due to associated bone marrow failure and malignancy risk.

\medterm{Rickets}-- See clinical_nutrition.

\medterm{Rotator Cuff Tears}-- Rotator cuff tears: partial (articular-sided PASTA, bursal-sided, intrasubstance) vs full-thickness. Ellman classification for PASTA: Grade I (<3mm), II (3-6mm), III (>6mm). Cofield sizing: small (<1cm), medium (1-3cm), large (3-5cm), massive (>5cm). Clinical: Neer/Hawkins impingement, drop arm, lag signs (external rotation lag, internal rotation lag), weakness. MRI: tear size, retraction, fatty infiltration (Goutallier classification: 0-4), muscle atrophy. Treatment: partial tears <50\%: nonopernonoperative (PT, corticosteroid injection, NSAIDs); full-thickness: repair (arthroscopic primary repair: single-row vs double-row, suture bridge, transosseous equivalent). Massive tears: nonop vs repair vs tendon transfer (latissimus dorsi for posterosuperior, pectoralis major for subscapularis) vs RTSA (reverse total shoulder arthroplasty, for massive retracted tears with rotator cuff arthropathy).

\medterm{Sacroiliac Joint}-- See anatomy.

\medterm{Scaphoid Fractures}-- Scaphoid fractures: most common carpal fracture (young males, FOOSH). Classification: Herbert: A (incomplete), B (complete: B1 distal waist, B2 proximal waist, B3 proximal pole, B4 distal tubercle, B5 transverse). Vascularity: retrograde from distal radius -> proximal pole at risk for AVN. Clinical: snuffbox tenderness, scaphoid compression test, pain with axial thumb loading. Imaging: PA/lateral/oblique/scaphoid views; MRI/CT if occult (high suspicion, negative XR). Treatment: distal pole/waistdistal pole/waist fractures: closed reduction and cast immobilization (below elbow thumb spica or scaphoid cast for 6-12 weeks). Proximal pole fractures: ORIF (headless compression screw) due to high risk of nonunion and AVN. Nonunion: bone graft with or without fixation.

\medterm{Scoliosis}-- A lateral spinal curvature >10 degrees on standing radiograph (Cobb angle). Idiopathic scoliosis (most common, 80\%): adolescent onset (\(\ge\)10 years), more common in girls. Congenital: due to vertebral anomalies. Neuromuscular: associated with cerebral palsy, muscular dystrophy. Degenerative: from asymmetric disc/ facet degeneration in older adults. Screening: Adam forward bend test (rib hump). Management: observation for mild curves (<25°), bracing for moderate (25-45°) in skeletally immature, surgical correction for severe (>45-50°) or progressive curves.

\medterm{Shoulder Instability}-- Anterior instability (95\%): Bankart lesion (anterior-inferior labral tear), Hill-Sachs lesion (posterolateral humeral head impression fracture), ALPSA (anterior labroligamentous periosteal sleeve avulsion), HAGL (humeral avulsion of glenohumeral ligament), GLAD (glenoid labral articular disruption). Recurrence: <20 years 80-90\%, >40 years <20\%. Treatment: first-time: immobilization (controversial, may not reduce recurrence), early PT; recurrent: arthroscopic Bankart repair (suture anchors), rempremplissage (Hill-Sachs remplissage), coracoid transfer (Latarjet, Bristow) for glenoid bone loss >20-25 percent, bone grafting (distal tibia allograft, iliac crest autograft).

\medterm{SLAP Tears}-- Superior labrum anterior-posterior (SLAP) tears: Snyder classification: I (fraying), II (detachment of biceps anchor, most common), III (bucket-handle tear of labrum, stable biceps), IV (bucket-handle extending into biceps). Clinical: O'Brien active compression test, Speed test, Yergason test, crank test. MRI arthrogram: sensitivity >90\%. Treatment: Type I: debridement; Type II: repair (young, overhead athletes), biceps tenodesis (older, >35-40, non-overhead); Type III: debridement; Type IV: debdebridement with biceps tenodesis or tenotomy; Type IV: repair plus biceps tenodesis/tenotomy.

\medterm{Soft Tissue Tumors}-- Lipoma: most common benign soft tissue tumor, subcutaneous, encapsulated. Liposarcoma: most common malignant STS in adults, deep, subtypes (well-differentiated/dedifferentiated, myxoid/round cell, pleomorphic). Synovial sarcoma: t(X;18), SS18-SSX, young adults, deep extremity, "triple sign" on MRI. Rhabdomyosarcoma: most common pediatric STS, embryonal (favorable), alveolar (PAX3/7-FOXO1, unfavorable). Malignant peripheral nerve sheath tumor (MPNST): NF1 association, painful, rapid growth. DesmoDesmoid tumor: locally aggressive, not metastatic, watch and wait, surgery (high recurrence), cryotherapy, NSAIDs, tyrosine kinase inhibitor. Myxoma: intramuscular, benign. Nodular fasciitis: reactive, self-limited. Elastofibroma dorsi: subscapular, benign. Imaging: MRI with contrast. Biopsy: core needle (preferred), incisional, excisional. Staging: AJCC, Enneking classification.

\medterm{Spinal Cord Injury}-- Spinal cord injury: ASIA impairment scale (A-E). Complete (ASIA A) vs incomplete. Syndromes: central cord (cervical hyperextension, upper > lower extremity weakness, burning hands), Brown-Sequard (hemisection: ipsilateral motor/proprioception loss, contralateral pain/temp loss), anterior cord (motor + pain/temp loss, proprioception spared), posterior cord (proprioception loss), conus medullaris (L1-L2, saddle anesthesia, bowel/bladder, mixed UMN/LMN), cauda equina (LMN, radicular pain, asymmetriasymmetric neurologic involvement). Treatment: high-dose methylprednisolone (controversial), blood pressure augmentation (MAP >85 mmHg for 7 days), surgical decompression and stabilisation. Timing: early (<24h) for cervical injury with incomplete deficit. Rehabilitation, bowel/bladder management, and psychosocial support.

\medterm{Spinal Stenosis}-- Narrowing of the spinal canal causing compression of the spinal cord or nerve roots. Lumbar stenosis (most common): neurogenic claudication (leg pain/weakness with walking, relieved by sitting/leaning forward), often bilateral. Cervical stenosis: myelopathy (hyperreflexia, Babinski, Hoffman, gait disturbance, hand clumsiness, bowel/bladder dysfunction). Risk factors: aging, degenerative disc disease, facet hypertrophy, ligamentum flavum hypertrophy, and congenital stenosis. Diagnosis: MRI. Treatment: physical therapy, NSAIDs, epidural steroid injections, and decompressive laminectomy with or without fusion for refractory or severe cases.

\medterm{Spine Disorders}-- The spine, or vertebral column, is a complex bony and soft-tissue structure that provides axial support, protects the spinal cord, and enables flexible movement of the trunk. It consists of 33 vertebrae divided into cervical, thoracic, lumbar, sacral, and coccygeal regions, separated by intervertebral discs that act as shock absorbers. Common spinal conditions include degenerative disc disease, herniated discs, spinal stenosis, scoliosis, and vertebral fractures, which can cause pain, neurological symptoms, and impaired mobility; management ranges from physical therapy and medications to surgical intervention.

\medterm{Spondylolisthesis}-- Forward slippage of one vertebral body relative to the vertebra below it. Isthmic (Type II, most common in children/adolescents): defect in pars interarticularis (spondylolysis), often L5-S1. Degenerative (Type III, most common in adults): facet joint and disc degeneration, often L4-L5. Presents with low back pain, radicular leg pain, and hamstring tightness. High-grade slips: may cause postural deformity and neurologic deficits. Diagnosis: standing lateral X-ray (Myerding grading: I-V based on percentage slip), MRI. Treatment: conservative (PT, bracing, activity modification) for low-grade. Surgical: reduction and fusion for high-grade, progressive slip, or persistent pain.

\medterm{Sprain}-- An injury to a ligament (connective tissue connecting bone to bone) caused by excessive stretching or tearing. Grading: Grade I (mild stretching, microscopic tears), Grade II (partial tear with laxity), Grade III (complete tear with instability). Common sites: ankle (lateral ligaments most common), knee (ACL, MCL), wrist, and thumb (gamekeeper's thumb). Presents with pain, swelling, bruising, and joint instability (Grades II-III). Diagnosis: clinical, stress testing, MRI for confirmation. Treatment: RICE (rest, ice, compression, elevation), NSAIDs, immobilization, physical therapy. Grade III injuries may require surgical reconstruction.

\medterm{Sprengel Deformity}-- A congenital elevation of the scapula due to failure of the scapula to descend from its embryonic cervical position to its normal thoracic position during weeks 9-12 of gestation. The scapula is higher, smaller, and rotated inferiorly compared to the normal side. Range of motion is limited, especially in shoulder abduction. Often associated with Klippel-Feil syndrome, scoliosis, and cervical ribs. An omovertebral bone (fibrous, cartilaginous, or bony connection between the scapula and cervical spine) is present in 30-50\% of cases. Physical exam shows an elevated scapula with a webbed-appearing neck on the affected side. X-ray confirms the elevated scapula and identifies associated cervical anomalies. Treatment: observation for mild cases; surgical correction (Woodward or Green procedure) for moderate to severe cases with functional or cosmetic concerns, ideally performed at 3-8 years of age.

\medterm{Strain}-- An injury to a muscle or tendon (connective tissue connecting muscle to bone) caused by excessive stretching, contraction, or overload. Grading: Grade I (mild stretching with few muscle fibers torn), Grade II (moderate partial tear), Grade III (complete muscle/tendon rupture). Common sites: hamstring, quadriceps, calf (gastrocnemius), lower back, and rotator cuff. Presents with acute pain, swelling, muscle spasm, and loss of function. Diagnosis: clinical, ultrasound or MRI for Grade II-III. Treatment: RICE (rest, ice, compression, elevation), NSAIDs, gradual return to activity with physical therapy. Grade III ruptures (e.g., Achilles tendon) may require surgical repair.

\medterm{Syndactyly}-- A congenital anomaly where two or more digits are fused together by skin (simple syndactyly) or by both skin and bone (complex syndactyly). Most common congenital hand anomaly, often bilateral. May be isolated (autosomal dominant, incomplete penetrance) or syndromic (Poland, Apert, Crouzon, Holt-Oram syndromes). The third web space is most commonly affected. Physical exam assesses the extent of fusion, nail involvement, and functional limitation. X-ray evaluates the degree of bony fusion and associated anomalies. Treatment: surgical separation (syndactyly release) performed at 6-18 months to prevent angular growth deformity. Z-plasties or skin grafts are used to reconstruct the web space. Separation of multiple digits requires staging to avoid vascular compromise.

\medterm{Tendinitis}-- Inflammation of a tendon, typically at its insertion (enthesis) or at the musculotendinous junction. Tendinitis is often caused by repetitive microtrauma/overuse but the term is increasingly replaced by 'tendinopathy' as histopathology shows tendinosis (collagen degeneration, neovascularisation, mucoid change) rather than acute inflammation. Common sites: rotator cuff (supraspinatus—shoulder impingement, subacromial pain), lateral epicondylitis (tennis elbow—extensor carpi radialis brevis origin), medial epicondylitis (golfer's elbow—common flexor origin), patellar tendinitis (jumper's knee—inferior pole of patella), Achilles tendinitis (insertional or mid-portion), de Quervain tenosynovitis (first dorsal compartment of the wrist—abductor pollicis longus, extensor pollicis brevis), and bicipital tendinitis (long head of biceps). Diagnosis: clinical (tenderness over the affected tendon, pain on resisted contraction, and passive stretch), ultrasound or MRI to confirm. Treatment: relative rest, activity modification, ice, NSAIDs, physiotherapy (eccentric exercises—proven effective for Achilles and patellar tendinopathy), corticosteroid injection (short-term relief, but may increase rupture risk—avoid for Achilles and patellar tendons), platelet-rich plasma (PRP) injection (evidence mixed), extracorporeal shock wave therapy, and surgery (tendon debridement, repair) for refractory cases.

\medterm{Tennis Elbow}-- Lateral epicondylitis, a tendinopathy of the common extensor tendon at the lateral epicondyle (extensor carpi radialis brevis). Caused by repetitive wrist extension and overuse, not limited to tennis. Presents with lateral elbow pain worsened by wrist extension and grip. Tenderness over the lateral epicondyle. Positive Cozen test and Mill test. Diagnosis: clinical; MRI or ultrasound to confirm. Treatment: rest, ice, activity modification, NSAIDs, physical therapy (eccentric exercises), and corticosteroid injection. Refractory cases: PRP injection, extracorporeal shock wave therapy, or surgical release.

\medterm{Thoracolumbar Spine Trauma}-- Thoracolumbar trauma: T10-L2 (thoracolumbar junction), L3-L5. AO classification: A (compression), B (distraction), C (translation/rotation). TLICS (Thoracolumbar Injury Classification and Severity): morphology (1-3), PLC (0-2), neurologic (0-3). <4 nonop, >4 op. Neurologic: complete (ASIA A) vs incomplete. Treatment: nonop (TLSO brace, early mobilization); op: posterior instrumentation (pedicle screws 1 above/1 below), decompression (laminectomy if retropulsed fragment, transpedicular or costotrcostotransversectomy for anterolateral compression). Post-op: TLSO brace, early mobilisation. Complications: infection, hardware failure, pseudarthrosis, proximal junctional kyphosis.

\medterm{Thumb CMC Arthritis}-- Thumb carpometacarpal (CMC) arthritis: Eaton-Littler classification: I (normal), II (slight narrowing, osteophytes <2mm), III (narrowing, osteophytes >2mm, subluxation), IV (pantrapezial arthritis). Clinical: base of thumb pain, pinch weakness, thenar wasting, squaring (subluxation). Treatment: I/II: splint, NSAIDs, corticosteroid injection, PT; III/IV: surgical: trapeziectomy (LRTI - ligament reconstruction tendon interposition with FCR/APL, suture button suspension, tightrope), total joint arttotal joint arthroplasty with hemi or total (trapezium-sparing, resurfacing).

\medterm{Tibial Shaft Fractures}-- Tibial shaft fractures: most common long bone fracture. AO/OTA: 42-A (simple), 42-B (wedge), 42-C (complex). Soft tissue critical: Gustilo for open. Treatment: IM nailing (gold standard, reamed, anteromedial or medial parapatellar entry, suprapatellar), plating (medial, anterolateral - for proximal/distal metaphyseal), external fixation (severe open, damage control). Compartment syndrome risk high (4 compartments). Complications: malunion (valgus/procurvatum), nonunion (exchange nail, plate + bobone grafting, exchange nailing, plate + bone graft).

\medterm{Torticollis (Wryneck)}-- A condition in which the head is tilted to one side (ear toward the shoulder) and rotated to the opposite side (chin toward opposite shoulder), due to contraction or shortening of the sternocleidomastoid muscle. Congenital muscular torticollis (most common): presents in infancy (2--3 weeks of age), caused by fibrosis of the sternocleidomastoid (possibly from intrauterine positioning, birth trauma, or compartment syndrome). A palpable neck mass (pseudotumour of infancy) may be present. Diagnosis: clinical, ultrasound to exclude other causes. Management: passive stretching exercises (95\% resolve with physiotherapy by 12 months). If persistent after 12--18 months: surgical release of the sternocleidomastoid. Acquired torticollis: causes—infection (retropharyngeal abscess, cervical lymphadenitis—Grisel syndrome, tonsillitis, mastoiditis—inflammatory torticollis), trauma (atlantoaxial rotatory subluxation/fixation—C1--C2 rotary fixation—most common in children after minor trauma or upper respiratory infection, presents with Cock Robin position), neurological (cervical dystonia/spasmodic torticollis—idiopathic, adult onset, treated with botulinum toxin injections), ocular (superior oblique palsy—head tilt to compensate for diplopia), and drug-induced (dystonic reaction to antipsychotics—acute, responsive to anticholinergics). Red flags: fever, meningism, neurological deficit, trauma, pain, or fixed deformity.

\medterm{Total Hip Arthroplasty}-- THA indications: end-stage OA, inflammatory arthritis, AVN, fracture, dysplasia. Approaches: posterior (most common, dislocation risk), direct anterior (DAA, muscle-sparing, fluoroscopy), anterolateral (Watson-Jones), lateral (Hardinge). Bearings: ceramic-on-poly (standard), ceramic-on-ceramic (young, squeaking risk), metal-on-poly (historical). Fixation: cemented (elderly, osteoporotic), cementless (porous coated, hydroxyapatite, young), hybrid (cemented stem + cementless cup). Complications: ddislocation (posterior approach 1-3 percent, DAA <1 percent), infection (1-2 percent), periprosthetic fracture, aseptic loosening, leg length discrepancy, nerve injury, venous thromboembolism.

\medterm{Total Knee Arthroplasty}-- TKA indications: end-stage OA, inflammatory arthritis, post-traumatic. Components: femoral (cobalt-chrome), tibial (titanium tray + polyethylene insert), patellar (polyethylene, resurfacing vs not). Alignment: mechanical (neutral), kinematic (patient-specific), restricted kinematic. Ligament balancing: measured resection, gap balancing. Patellar resurfacing: selective vs routine. Complications: stiffness (manipulation under anesthesia, arthroscopic lysis), instability (collateral, PCL), patellofpatellofemoral complications, aseptic loosening, periprosthetic fracture, infection (1-2 percent), venous thromboembolism, arthrofibrosis, and extensor mechanism disruption.

\medterm{Trigger Finger}-- Trigger finger (stenosing tenosynovitis): A1 pulley thickening, flexor tendon nodule -> catching/locking. Clinical: triggering, palpable nodule at distal palmar crease, morning stiffness. Grades: I (pain only), II (catching, active extension), III (catching, passive extension), IV (locked flexed). Treatment: corticosteroid injection (60-70\% cure, less in diabetics); percutaneous release (needle); open A1 pulley release (gold standard for refractory/recurrent). Thumb: MP flexion contracture risk.

\medterm{Whiplash Injury}-- An acceleration-deceleration injury to the cervical spine, most commonly resulting from rear-end motor vehicle collisions. Mechanism: sudden hyperextension followed by hyperflexion of the neck causes stretching or tearing of cervical muscles, ligaments, facet joint capsules, and annulus fibrosus of intervertebral discs. The S-shaped curve of the cervical spine during whiplash (initial straightening of the cervical lordosis followed by flexion at the lower segments and extension at the upper segments) produces shear stress on the facet joints. Clinical presentation: neck pain and stiffness (onset within 24 hours, may be delayed up to 72 hours), headache (typically occipital or suboccipital, cervicogenic origin), interscapular pain, upper limb pain and paraesthesiae (referred or radicular, from nerve root stretching), dizziness, fatigue, and temporomandibular joint pain. Quebec Task Force Classification: Grade 0 (no neck complaint, no physical sign), Grade I (neck pain, stiffness, tenderness only, no neurological signs), Grade II (neck complaint AND musculoskeletal signs — decreased range of motion, point tenderness), Grade III (neck complaint AND neurological signs — decreased or absent deep tendon reflexes, weakness, sensory deficits), Grade IV (neck complaint AND fracture or dislocation on imaging). Chronic whiplash syndrome: persistent pain and disability beyond 6 months, associated with high initial pain intensity, female sex, older age, and psychological factors (catastrophizing, fear avoidance). Diagnosis: clinical; imaging (cervical spine X-ray, CT, MRI) indicated if red flags present (fracture suspicion, neurological deficit, age >65, high-energy mechanism). Management: active treatment (early mobilisation, neck exercises, NSAIDs, muscle relaxants) is superior to immobilisation (soft collar). Cervical collars should be discontinued within 72 hours. Physiotherapy, manual therapy, and acupuncture for persistent pain. Prognosis: 50\% recover within 3 months, 75\% by 6 months. Chronic pain develops in 25\%.

\medterm{Wrist Drop}-- See Radial Nerve Palsy (Wrist Drop).
